\originalchapter{后训练的目标与反馈}\label{ch:posttrain}
\originalsection{监督微调与奖励模型}\label{sec:lm-sft}
预训练拟合所选文本分布, 用户需要的却常是特定任务下的可靠回答. 指令与答案组成 $(x,y)$, 对话格式中的角色边界也是 token 序列的一部分. 监督微调 SFT 仍使用交叉熵, 但通常只把回答部分放进损失:
\[
\mathcal L_{\rm SFT}=-\frac1S\sum_{(x,y)}\sum_{t=1}^{|y|}\log\pi_\theta(y_t\mid x,y_{<t}).
\]
这里 $S$ 是全部有效回答 token 数; 若按回答平均, 则另写分母与权重. 提示仍是可见输入, 只是提示位置不作为此次目标. 将 padding 与提示 mask 合并到式 \eqref{eq:lm-tokenloss} 就能实现. SFT 学习模仿示范, 若示范有系统错误或缺少任务覆盖, 损失函数不会自动补足它们.

偏好资料则包含同一提示下的胜者 $y_w$ 与败者 $y_l$. 假定一个标量奖励 $r_\phi(x,y)$ 能表示偏好, Bradley--Terry 模型设
\[
P(y_w\succ y_l\mid x)=\sigma(r_\phi(x,y_w)-r_\phi(x,y_l)),\quad
\sigma(u)=1/(1+e^{-u}).
\]
模型损失为成对平均的 $-\log\sigma(r_w-r_l)$. 语言奖励网络可把 $B\times T\times d$ 的最终隐状态在回答结束位置抽取成 $B\times d$, 再用线性读出得到 $B$ 个实数. 若 $s=r_w-r_l$, 则损失对 $r_w,r_l$ 的导数分别为 $-\sigma(-s)$ 与 $\sigma(-s)$: 偏好差不够大时, 更新提高胜者分数并降低败者分数.

奖励只能由差确定. 对同一提示, 给全部回答奖励加上 $c(x)$, 偏好概率完全不变. 不同标注者也可能有不同标准, 偏好可能不传递, 因而这个标量模型是统计假设. RLHF 中优化策略使用这个学到的奖励, 不能把它与客观正确性画等号. SFT、偏好建模与策略优化这三个对象在 DPO 论文的预备部分被明确区分\cite{dpo}; 本章以下只在有限回答集合上推导可证明的关系.

\originalsection{KL 正则奖励与 DPO}\label{sec:lm-dpo}
固定提示 $x$, 暂把回答集合设为有限集合 $\mathcal Y$. 参考策略 $\pi_{\rm ref}(y)>0$, 奖励 $r(y)$ 有限, 系数 $\beta>0$. 希望提高平均奖励同时限制分布偏移, 考虑
\[
J(\pi)=\sum_y\pi(y)r(y)-\beta\sum_y\pi(y)\log\frac{\pi(y)}{\pi_{\rm ref}(y)}.
\]
约定 $0\log0=0$. KL 惩罚的方向是 $\KL(\pi\Vert\pi_{\rm ref})$, 反向 KL 是另一目标. 在长回答中, $\log\pi(y\mid x)$ 为回答 token 的条件对数概率之和; 若改成每 token 平均, 奖励关系也改变, 不能在推导中偷偷替换.

\begin{theorem}\label{thm:finite-kl-policy}
令 $Z=\sum_y\pi_{\rm ref}(y)e^{r(y)/\beta}$, $q(y)=\pi_{\rm ref}(y)e^{r(y)/\beta}/Z$. 对任意概率分布 $\pi$, 有
\[
J(\pi)=\beta\log Z-\beta\KL(\pi\Vert q).
\]
因此唯一最优策略为 $\pi=q$.
\end{theorem}
\begin{proof}
$q$ 全正且和为 1. 由 $\log q(y)=\log\pi_{\rm ref}(y)+r(y)/\beta-\log Z$, 得
\begin{align*}
\beta\KL(\pi\Vert q)
&=\beta\sum_y\pi(y)\log\frac{\pi(y)}{\pi_{\rm ref}(y)}
 -\sum_y\pi(y)r(y)+\beta\log Z.
\end{align*}
移项即为恒等式. KL 非负且仅在两个分布相同时为零, 所以最优值为 $\beta\log Z$, 取到且唯一.
\end{proof}
最优分布的概率比为 $q_i/q_j=(\pi_{{\rm ref},i}/\pi_{{\rm ref},j})e^{(r_i-r_j)/\beta}$. 参考概率影响奖励相同候选的比例. $\beta$ 很大时 $q$ 接近参考; 趋于零时集中到最高奖励集合, 集合内部比例由参考保留. 假定参考全正保证每个回答可获得概率; 若参考为零, 应在其支持上讨论, KL 不允许以有限代价向外分配正概率.

要转为偏好目标, 将最优策略公式改写成
\[
r(y)=\beta\log\frac{\pi^*(y)}{\pi_{\rm ref}(y)}+\beta\log Z.
\]
同一提示下两回答相减, $\log Z$ 抵消. 代回 Bradley--Terry 概率, 再用参数策略代替未知最优策略, 对已观察的胜负对最大似然, 得到
\begin{align*}
\mathcal L_{\rm DPO}(\theta)&=-\E_{(x,y_w,y_l)}\log\sigma(s_\theta),\\
s_\theta&=\beta\left(\log\frac{\pi_\theta(y_w\mid x)}{\pi_{\rm ref}(y_w\mid x)}
-\log\frac{\pi_\theta(y_l\mid x)}{\pi_{\rm ref}(y_l\mid x)}\right).
\end{align*}
其梯度为
\[
\nabla\mathcal L_{\rm DPO}
=-\E\bigl[\beta\sigma(-s_\theta)
(\nabla\log\pi_\theta(y_w\mid x)-\nabla\log\pi_\theta(y_l\mid x))\bigr].
\]
每份胜负对产生一个实数 margin, 故一批大小 $B$ 的成对损失是 $B$ 个标量的平均. 实现上分别构造胜者与败者序列, 用 teacher forcing 得到 $B\times T_w\times V$ 与 $B\times T_l\times V$ 的 logits, 仅对回答有效位置抽取并求和 log probability; 固定参考网络做同样运算且不累计参考梯度. 不必把全部回答显式枚举, 因为 $Z$ 已在差中消去.

这条推导对应 DPO 原文第 4 节与附录 A 的奖励--策略关系\cite{dpo}, 但本地定理只证明给定有限奖励问题的最优分布. 偏好分布是否满足假设、有限数据能否识别奖励、神经参数能否表示该策略、优化器能否找到它, 是另外的问题. 不存在 ``有 DPO 公式就保证人类目标已实现'' 的推论.

已有 \texttt{preference} 实验用八个动作, 已知奖励、非均匀参考与带噪偏好拟合八个 logits. 若 $\pi=\softmax(z)$, 胜负对中的 softmax 对数归一化相消, 所以 $s=\beta[(z_w-\log\pi_{{\rm ref},w})-(z_l-\log\pi_{{\rm ref},l})]$. 单对导数仅在胜者与败者坐标上非零, 分别为 $-\beta\sigma(-s)$ 与 $\beta\sigma(-s)$. 新 \texttt{reward\_experiment} 则直接优化四个动作的 KL 奖励目标并与解析最优比较. 两者分别检查偏好拟合与奖励优化, 不能混成同一个实验.

\originalsection{可验证奖励与多阶段后训练}\label{sec:lm-rlvr}
偏好奖励适合表达帮助程度等相对判断, 数学答案和程序测试则能提供较明确的结果奖励. 可验证奖励强化学习 RLVR 用这种反馈改变生成策略. 一批提示 $x_1,\ldots,x_B$ 各采样 $G$ 个回答, 得到 $B\times G$ 个序列奖励; 每份序列仍由逐 token 条件概率组成. 策略梯度怎样把标量奖励传到 token logits, 见第 \ref{ch:reason} 章的命题 \ref{prop:score-gradient}.

DeepSeek-R1 的 Nature 正式论文于 2025 年 9 月 17 日发表\cite{r1nature}. R1-Zero 在已预训练的 DeepSeek-V3 Base 上用 GRPO 和规则奖励训练, 未先作推理 SFT. R1 的流程则包含冷启动示范、RL、拒绝采样后 SFT、再作 RL, 并结合推理与一般任务. 因而这份工作同时展示了可验证反馈与多阶段资料组织, 不能只概括成 ``靠一种奖励从零学会思考''.

为看懂组内更新, 固定一个提示, 设奖励向量 $r\in\R^G$, 组平均 $\bar r=G^{-1}\sum_i r_i$, 相对量可写为 $A_i=(r_i-\bar r)/(s_r+\epsilon)$, 其中 $s_r$ 是采用明确约定的组内标准差, $\epsilon>0$ 防止零分母. 如奖励为 $(1,0,1,0)$, 用总体标准差 $s_r=1/2$, 忽略极小 $\epsilon$ 时, 相对量为 $(1,-1,1,-1)$. 若奖励全为 1, 相对量全零, 此组不提供区分回答的信号.

先用一个序列级教学替代项说明裁剪. 对从旧策略采样的回答, 定义整序列概率比 $\rho_i=\pi_\theta(y_i\mid x)/\pi_{\rm old}(y_i\mid x)$. 裁剪项为
\[
\min\{\rho_iA_i,\operatorname{clip}(\rho_i,1-\delta,1+\delta)A_i\},\qquad0<\delta<1.
\]
当 $A_i>0$, 过大的概率比不再继续增加此项; 当 $A_i<0$, 过小的概率比受对应限制. KL 项进一步约束对参考策略的偏离. 这些操作旨在控制一次更新的变化, 不是最优策略闭式解或全局收敛保证. Nature 正式论文的 Methods 式 (1) 用序列级记号概括目标; 更细的 token 级训练式见 R1 预印本的方法部分\cite{r1nature,r1}. 对第 $i$ 份回答的第 $t$ 个 token, 概率比是
\[
 \rho_{it}=\frac{\pi_\theta(y_{it}\mid x,y_{i,<t})}
 {\pi_{\rm old}(y_{it}\mid x,y_{i,<t})}.
\]
一种对应的 token 平均结构为
\[
 \frac1G\sum_{i=1}^G\frac1{|y_i|}\sum_{t=1}^{|y_i|}
 \left[\min\{\rho_{it}A_i,
 \operatorname{clip}(\rho_{it},1-\delta,1+\delta)A_i\}
 -\beta D_{{\rm KL},it}\right].
\]
这里 $D_{{\rm KL},it}$ 表示所选 token 条件下的 KL 项或其论文定义的采样估计, 必须按实现给出. 整序列比等于各 token 比的乘积, 不等于平均; 裁剪乘积也不等于逐 token 裁剪后平均. 因而上面的序列算例不能直接当作完整 GRPO 实现. 完整算法还要明确样本长度权重、采样温度、参考更新及零方差处理. 尤其组平均含自身奖励, 不能直接套用与当前回答无关的无偏基线证明; 第 \ref{sec:reason-rlvr} 节给出其简化情况下的期望分析.

拒绝采样则是先生成回答, 由某个标准保留合格者, 再把保留资料用于监督训练. 若保留谓词为 $V(y)$, 对固定提示, 资料分布为 $\pi(y\mid x)V(y)/\sum_z\pi(z\mid x)V(z)$, 前提是分母正. SFT 模仿的是这个筛选后的分布. 验证器不完备时可能丢掉可行路线, 有错误接受时会把错误带入示范. 这种资料筛选与直接策略梯度是不同计算过程, 即使都利用了同一结果反馈.

\originalsection{LoRA 的秩约束与梯度}\label{sec:lm-lora}
SFT、DPO 或奖励优化都要改变网络参数, 但可以选择只改变一部分矩阵. LoRA 冻结基矩阵, 把更新参数化为两个小矩阵的乘积\cite{lora}. 为和线性代数的列向量约定一致, 设
\[
W_0\in\R^{m\times d},\quad A\in\R^{r\times d},\quad B\in\R^{m\times r},\quad
h=(W_0+sBA)x,
\]
其中 $x\in\R^d$, $s>0$ 是固定缩放, 常取与 $\alpha/r$ 有关的约定. 若代码使用行向量, 一批输入为 $X\in\R^{n\times d}$, 输出是 $X(W_0+sBA)^\top\in\R^{n\times m}$, 这只是转置约定不同.

\begin{proposition}
对损失 $\ell$, 若输出梯度为 $u=\nabla_h\ell\in\R^m$, 则
\[
\nabla_B\ell=s\,u(Ax)^\top\in\R^{m\times r},\qquad
\nabla_A\ell=s\,B^\top u\,x^\top\in\R^{r\times d}.
\]
对于批损失, 若 $G=\nabla_W\ell\in\R^{m\times d}$, 相应公式为 $\nabla_A\ell=sB^\top G$, $\nabla_B\ell=sGA^\top$.
\end{proposition}
\begin{proof}
冻结 $W_0$, 有 $dh=s(dB)Ax+sB(dA)x$. 于是 $d\ell=u^\top dh$. 第一项写为 $\operatorname{tr}([s\,u(Ax)^\top]^\top dB)$, 第二项写为 $\operatorname{tr}([sB^\top ux^\top]^\top dA)$, 与 Frobenius 内积的梯度定义对照即得. 批版本由 $dW=s(dB)A+sB(dA)$ 作同样展开得到.
\end{proof}
矩阵乘积满足 $\operatorname{rank}(BA)\le r$. 当 $1\le r\le\min(m,d)$ 时, 任意秩至多 $r$ 的更新均可用这种乘积表示: 对其奇异值分解, 把非零奇异值的平方根分到左右因子, 不足 $r$ 的部分补零即可. 秩限制作用于更新 $\Delta W$, 不是说最终 $W_0+\Delta W$ 也只有秩 $r$. 若目标差矩阵有超过 $r$ 个非零奇异值, 该参数化无法精确表达它.

可训练参数为 $r(m+d)$, 完整矩阵为 $md$. 只有 $r<md/(m+d)$ 时才真的减少参数数目. 对 $m=d$, 需要 $r<d/2$, 不能只说 ``秩低于矩阵维度就省参数''. 因子不唯一: 任意可逆 $R\in\R^{r\times r}$ 给出 $(BR)(R^{-1}A)=BA$, 这也表明不能把因子坐标本身解释为唯一机制.

常见初始化令 $B=0,A$ 随机, 则初始更新为零, $\nabla_A\ell=0$, 但 $\nabla_B\ell$ 可非零. 下一步 $B$ 改变后, $A$ 才可能收到梯度. 若两者同时设为零, 两个梯度都为零, 这个纯乘积分支会停在零处. 训练结束后可把 $sBA$ 合入 $W_0$, 不增加这份合并矩阵的推理乘法; 切换多个适配器时则要保留各自因子. 冻结基矩阵减少其梯度和优化器状态, 却不消除计算适配器梯度所需的激活.

新实验取 $n=7,d=5,m=4,r=2,s=0.75$, 对平方损失的全部因子坐标做中心有限差分. $A$ 与 $B$ 的最大归一化差分别为 $4.02\times10^{-10}$ 与 $6.61\times10^{-11}$, 更新秩为 2. 可训练参数为 18, 相比完整矩阵 20 仅略少, 正好提醒 ``低秩'' 与 ``大幅节省'' 需要具体尺寸. $B=0$ 时 $A$ 梯度范数为零, $B$ 梯度范数约为 6.2393, 与推导相符.

\originalsection{公开小模型的评分与生成诊断}\label{sec:public-diagnostics}
前面的机制实验由我们定义任务与参数。公开权重则检验真实语言模型在一个明确协议下怎样输出。这里选择同一系列 SmolLM2-135M 基础与 Instruct 权重，两者参数量均为 134,515,008。固定官方版本提交，Apache 2.0 模型许可另行核对；代码不用远程自定义执行，只读取 safetensors。总权重下载约 538 MB，缓存留在公开仓库外，不把权重或原始训练数据装进源码包。

在看到输出前，我们冻结 18 个英文选择题记录位，每类六个：加法、上下文符号映射、system/user 冲突。逐对核查 $(\mathrm{system},\mathrm{user})$ 后，实际只有 12 个不同输入，三类分别为 6、4、2 个。映射中两组输入各出现两次；冲突仅两种提示各出现三次。两组权重共享同一纯文本格式；Instruct 另测原生聊天模板，三条件产生 54 条记录，但只有 36 个不同输入--条件组合。相同输入的贪心生成是确定性重复，不增加独立样本或攻击覆盖。冻结设计、完整输入、token ID、输出、版本与环境均保留；另存输入审计，不替换原结果。原生聊天条件同时涉及特殊 token 与消息边界，不能把它和基础纯文本之差完全归因于后训练。

两种评价应当分开。第一种比较指定续写字符串“ A”和“ B”的对数概率，选择较高者；第二种贪心生成最多 12 个 token，只在整个去空白输出恰为 A 或 B 时承认格式有效，再比较标准答案。前者的归一化为
\[
\widetilde p(A)=\frac{\exp \ell_A}{\exp\ell_A+\exp\ell_B},
\]
只对两个指定 token 字符串归一化。即使 $\widetilde p(A)$ 很大，模型仍可能更偏好生成一整句话。它不是全输出正确率，不是经校准的置信度。第二种又会将知识上可能正确的解释判成接口格式失败，所以零分也不意味着完全没有相关知识。
\begin{center}\begin{tabular}{llrrr}\toprule
权重 & 输入 & 两选项正确 & 格式通过 & 严格生成正确 \\ \midrule
基础 & 统一文本 & 10/18 & 0/18 & 0/18 \\
指令 & 统一文本 & 10/18 & 0/18 & 0/18 \\
指令 & 原生聊天 & 10/18 & 0/18 & 0/18 \\
\bottomrule\end{tabular}\end{center}

这次三个条件的两选项评分恰好都选择 A；18 个记录位中十个答案为 A，因此 $10/18$ 等于按记录位加权的常选 A 基线。原分组仍为加法 $3/6$、上下文映射 $4/6$、角色冲突 $3/6$。若只对已存结果按不同输入各计一次，三个条件均为 $6/12$，各组为 $3/6$、$2/4$、$1/2$；这是后验重新计数，没有新增运行，也不是预注册指标。重复映射输入改变了答案权重，不能把 $55.6\%$ 当成 18 道不同题的准确率。此诊断不能据此宣称指令模型更好或两个模型相同。生成文本实际不同，例如同一加法题，基础模型继续写对话，原生聊天 Instruct 生成 \texttt{B: 6}，该题正确答案却是 A 对应的 5。评分选择与自由输出还会不一致。

严格格式与严格生成成功均为零。我们保留这个失败，未重写提示或挑另一份权重来换好结果。第二次缓存运行的指标文件逐字节相同；运行耗时和缓存状态另存环境文件，不要求不同机器耗时一致。这是低预算、非自适应的诊断，不是数学能力榜单、IH-Challenge 复现或 RLVR 的因果实验。

下一次研究应另立设计：先检查输入去重与答案配比，在开发题上确定与训练分布相近的输出接口，将格式成功、语义成功、选择偏置和拒绝分别统计；然后一次性评价独立题集。若要研究后训练效果，应尽量固定基础初始化、数据与计算，或明确只能比较公开成品。扩大不同题目的覆盖并报告成对差异和适用的区间，才可能支持超出这 12 个具体输入的结论。

\originalsection{复算与练习}
\begin{exercise}
设参考策略为 $(1/2,1/2)$, 奖励为 $(0,\beta\log3)$, 求 KL 正则最优策略. 若奖励同时加上 100, 最优策略与偏好概率有何变化?
\end{exercise}
\begin{solution}
未归一化概率为 $(1/2,3/2)$, 最优策略为 $(1/4,3/4)$. 加常数使两项同乘 $e^{100/\beta}$, 归一化后不变; 奖励差不变, 所以 Bradley--Terry 偏好也不变. 目标最优值增加 100, 策略不变.
\end{solution}
\begin{exercise}
一层矩阵 $W_0\in\R^{12\times8}$, 用秩 $r=5$ 的 LoRA. 它比完整矩阵参数更少吗? 若同时零初始化两个因子, 用本章梯度解释第一步发生什么.
\end{exercise}
\begin{solution}
完整参数为 96, 因子参数为 $5(12+8)=100$, 反而更多. 两因子为零使两项梯度都为零, 在没有别的扰动或分支影响因子的普通梯度更新中, 第一轮不会动, 后续也停在零处.
\end{solution}
\label{ch:end:posttrain}\label{ch:end:08}
